\section{Proper forcing and master conditions}

\begin{annexitem}{annex:p06.pfm.hchi}{thm:p06-pfm-hchi}{Hereditary size and the set bound}
\begin{proof}
All constructions in this proof are in \(\mathsf{ZFC}\).
First form \(T=\operatorname{tc}(\{x\})\): start with \(T_0=\{x\}\),
put \(T_{n+1}=T_n\cup\bigcup T_n\), and take \(\bigcup_{n<\omega}T_n\).
Ordinary set recursion and Replacement give this sequence and its union.
If \(u\in v\in T\), then \(v\in T_n\) for some \(n\), so
\(u\in T_{n+1}\). Induction on \(n\) proves minimality among transitive
sets containing \(x\). The same construction for each \(y\in T\)
gives \(T_y=\operatorname{tc}(\{y\})\subseteq T\).

We give the rank bound rather than assume that this hereditary-size class
is a set. An approximation is an ordinal-valued function on a transitive
subset \(S\subseteq T\) satisfying
\[
 f(y)=\sup\{f(z)+1:z\in y\}\qquad(y\in S).
\]
Two approximations agree on their common domain. Otherwise Foundation
applied to the set of disagreement points gives a point \(y\) with no
member at which they disagree; both domains contain every \(z\in y\),
and the two displayed suprema are identical, a contradiction.

For each \(y\in T\) there is exactly one approximation on \(T_y\).
Uniqueness was just proved. For existence use Foundation induction on \(T\).
If approximations on \(T_z\) exist for all \(z\in y\), their uniqueness
makes \(z\mapsto f_z\) a functional formula. Replacement on the set \(y\)
and Union give their common compatible graph \(g\), whose domain is
\(\bigcup_{z\in y}T_z\). If this domain already contains \(y\), restrict
\(g\) to \(T_y\). Otherwise extend it at \(y\) by the displayed supremum
of its known values at the members of \(y\). Older equations are unchanged;
the new one holds by construction. Its domain is \(T_y\).
Thus the induction step supplies the required approximation. Formally,
if any \(y\) failed, Foundation on the set of failures would give a failure
whose members all satisfy this step, a contradiction.
Replacement on \(T\) now collects these unique graphs; their union is one
function \(r:T\to\operatorname{Ord}\) with the displayed equations.

The set \(R=r[T]\) is an ordinal. To see that it is downward closed, use
Foundation induction: if \(\beta<r(y)\), the supremum equation supplies
\(z\in y\) with \(\beta\leq r(z)\). Equality places \(\beta\) in \(R\);
strict inequality uses the induction assertion for \(z\).
All members of \(R\) are ordinals, so downward closure makes \(R\) an
ordinal by ordinal comparison and Foundation.
If \(|T|<\chi\), then \(|R|\leq|T|<\chi\), using a least preimage in a
well-order of \(T\). Since \(\chi\) is an initial ordinal, \(R<\chi\).
Foundation induction and the defining \(V\)-hierarchy clauses show
\(y\in V_{r(y)+1}\): all members of \(y\) lie in \(V_{r(y)}\), so
\(y\in\mathcal P(V_{r(y)})\).
Here the existence/clauses of the hierarchy and definable set recursion are
the completed \cref{thm:p06-hierarchy-recursion-existence,thm:p06-definable-ordinal-recursion},
not a rank sketch. An infinite cardinal is a limit ordinal, so
\(r(x)+1<\chi\) and \(x\in V_\chi\).
Separation on the set \(V_\chi\), with the hereditary-size predicate,
therefore forms exactly \(H_\chi\).

If \(y\in x\in H_\chi\), then \(T_y\subseteq T_x\), hence \(y\in H_\chi\).
For \(b\subseteq a\in H_\chi\), its transitive closure is contained in
\(\{b\}\cup T_a\); thus \(b\in H_\chi\).
The transitive closure of a finite tuple is bounded by the finitely many
input transitive closures and finitely many levels of pair coding.
Finite sums/products of cardinals below the infinite \(\chi\) stay below
\(\chi\), so these tuples also belong to \(H_\chi\).

For nonempty countable \(a\), choose a surjection \(f:\omega\to a\).
Its graph has countably many pairs; its transitive closure consists of
their finitely coded components together with subsets of \(T_a\) and
\(\omega\). Its size is at most
\(\max(\aleph_0,|T_a|)<\chi\), so \(f\in H_\chi\).
Finally \(|\operatorname{tc}(\{\omega_1\})|=\aleph_1<\chi\), and the
declared \(\kappa\) bounds the poset datum's transitive closure.
None of these arguments asserts that the collection of \emph{all}
subsets of a member is itself a member of \(H_\chi\).
\end{proof}
\end{annexitem}

\begin{annexitem}{annex:p06.pfm.hulls}{thm:p06-pfm-hulls}{Witness closure and model traces}
\begin{proof}
The membership language has countably many finite formulas.
At each stage the tuples from a countable set and the indices of the
countably many witness functions can be enumerated by finite natural-number
tuple codes. Their image is countable. Starting with an enumeration of the
seed, this supplies stage enumerations and then a single enumeration of
all pairs consisting of a stage and an entry. Thus the union hull is
countable. Every finite tuple in the union lies in one stage; its witness
lies in the next stage. The Tarski--Vaught test
\cref{thm:tarski-vaught} applies to every existential membership formula
and proves elementarity. Adding countably many specified operations
does not remove these witnesses.

Let \(C\) be the countable sets containing the named constants and closed
under all the specified operations. Every countable seed has a hull in
\(C\) containing it, so \(C\) is unbounded. In the union of an increasing
countable chain, every finite tuple occurs in one member; its operation
value occurs there as well. The union is still countable and closed,
so \(C\) is club.
For arbitrary clubs \(C_0,C_1\), repeatedly choose a member of \(C_0\)
containing the seed, a member of \(C_1\) containing that member, and then
alternate. The union belongs to both clubs by their chain closure.
This proves unboundedness of the intersection; chain closure of the
intersection follows separately from that of each club.

The empty set is uniquely definable in the transitive \(H_\chi\),
so \(0\in N\). Unique successor formation and finite induction in the
metatheory give every natural number in \(N\).
The least infinite ordinal computed in \(H_\chi\) is the real \(\omega\):
all its finite ordinals and \(\omega\) belong to \(H_\chi\), and the
membership order and finiteness of these ordinals are unchanged.
Consequently \(\omega\in N\).

Let \(a\in N\) be nonempty countable. The preceding hereditary-size lemma
places a surjection \(\omega\to a\) in \(H_\chi\).
Its existence is therefore true in \(H_\chi\); elementarity supplies one
\(f\in N\). Transitivity makes the function/domain/surjectivity assertions
about this \(f,\omega,a\) correct in the metatheory.
For each \(n\in\omega\subseteq N\), the unique-value formula for \(f(n)\)
has a witness in \(N\), and that witness is the actual value.
Surjectivity proves \(a\subseteq N\). The empty case is immediate.

The least uncountable ordinal of \(H_\chi\) is the real \(\omega_1\).
For every \(\beta<\omega_1\) its countability witness belongs to \(H_\chi\)
by the same graph bound, while no surjection \(\omega\to\omega_1\)
exists there because none exists in the metatheory.
Thus \(\omega_1\in N\) by its unique definition.
If \(\beta\in N\cap\omega_1\), it is countable, so \(\beta\subseteq N\).
Therefore \(N\cap\omega_1\) is an ordinal. It is countable because \(N\)
is countable, hence below \(\omega_1\).
Successor formation puts \(\beta+1\) in this intersection for every
\(\beta\) already in it; it has no greatest member and contains \(\omega\).
It is the asserted countable limit ordinal.
\end{proof}
\end{annexitem}

\begin{annexitem}{annex:p06.pfm.predensity}{thm:p06-pfm-predensity}{The extension quantifier and generic filters}
\begin{proof}
If \(r\leq_P q\) is compatible with \(x\in X\), choose a common
strengthening \(s\leq_P r,x\). Then \(s\in X^\downarrow\).
Conversely a strengthening \(s\leq_P r\) in \(X^\downarrow\) supplies
\(x\in X\) with \(s\leq_P x\). This proves the equivalence.

For dense \(D\), \(D^\downarrow\) is dense and open.
It is a subset of \(P\), hence belongs to \(H_\chi\), and is uniquely
defined there from \(D,\mathbf P\); when these belong to \(N\), it
belongs to \(N\). If only dense open sets are tested, take
\(y\in D^\downarrow\cap N\) compatible with a prescribed \(r\leq_P q\).
Elementarity with parameters \(y,D,\mathbf P\) supplies
\(d\in D\cap N\) with \(y\leq_P d\).
A common strengthening of \(r,y\) is also below \(d\), proving the
full dense-set test. The reverse implication is restriction of the test.

For the filter assertion form, in the ground model,
\[
 E=X^\downarrow\cup\{s\in P:s\perp_P q\}.
\]
It is dense: a condition incompatible with \(q\) already belongs to it;
otherwise strengthen to a common extension with \(q\) and use predensity.
A generic \(G\) meets \(E\). If \(q\in G\), directedness prevents any
member of \(G\) being incompatible with \(q\).
Thus some \(s\in G\cap X^\downarrow\) lies below an \(x\in X\),
and upward closure of \(G\) puts \(x\) in \(G\).
If predensity fails, its failed extension \(r\) cannot coexist in any
directed filter with any \(x\in X\). This is a concrete obstruction,
not an existence claim for generic filters over the current universe.
\end{proof}
\end{annexitem}

\begin{annexitem}{annex:p06.pfm.antichains}{thm:p06-pfm-antichains}{A maximal antichain chosen inside the model}
\begin{proof}
Well-order \(D\) using Choice. Recursively retain an element of \(D\)
exactly when it is incompatible with every previously retained element.
Definable recursion on that set-sized order supplies a set \(A\subseteq D\).
It is an antichain. Every skipped member was compatible with an earlier
retained member; every retained member is compatible with itself.
Thus every \(d\in D\) is compatible with \(A\).
For \(r\in P\), density supplies \(d\leq_P r\) in \(D\), and a common
strengthening of \(d\) and its compatible member of \(A\) is also below
\(r\). Hence \(A\) is maximal in \(P\).

All subsets of \(P\) belong to \(H_\chi\). The assertion that such an
\(A\subseteq D\) is an antichain compatible with every \(r\in P\)
uses only bounded quantification over these sets and the order relation.
It is true in \(H_\chi\). Elementarity supplies an \(A\in N\) with
exactly that real property whenever \(D,\mathbf P\in N\).

If \(q\) passes the maximal-antichain test, this \(A\subseteq D\)
gives the dense-set test.
Conversely, for a maximal \(A\in N\), the set \(A^\downarrow\)
is dense open and belongs to \(N\).
The dense-set test gives \(y\in A^\downarrow\cap N\) compatible with
each prescribed extension of \(q\).
Elementarity supplies \(a\in A\cap N\) above \(y\), and compatibility
transfers to \(a\). This proves the other direction.
\end{proof}
\end{annexitem}

\begin{annexitem}{annex:p06.pfm.ccc}{thm:p06-pfm-ccc}{Why ccc makes every condition a master condition}
\begin{proof}
In the countable case \(P\in N\) and the countable-element lemma gives
\(P\subseteq N\). For dense \(D\in N\) and \(r\leq_P q\), choose
\(d\leq_P r\) in \(D\). It belongs to \(N\) because every condition does.
Thus \(D\cap N\) is predense below every \(q\).

For ccc, choose a maximal antichain \(A\subseteq D\) in \(N\)
by the preceding construction and elementarity.
It is countable in the metatheory, so \(A\subseteq N\).
Maximality makes every \(r\in P\) compatible with an \(a\in A\),
which belongs to \(D\cap N\). This again works below every \(q\),
whether or not \(q\in N\). Taking \(q=p\) for the condition in the
properness test proves properness at every admissible height.
No assertion that all of an arbitrary dense \(D\in N\) lies in \(N\)
was needed.
\end{proof}
\end{annexitem}

\begin{annexitem}{annex:p06.pfm.closed}{thm:p06-pfm-closed}{A lower bound for the model's dense-set sequence}
\begin{proof}
The family \(\{D\in N:D\subseteq P\text{ dense}\}\) is countable.
It is nonempty since \(P\in N\), so enumerate it as \(D_0,D_1,\ldots\),
with repetitions if needed. This is an enumeration in the metatheory;
no enumeration inside \(N\) is assumed.
Put \(p_0=p\). If \(p_n\in P\cap N\), density of \(D_n\) and elementarity
give \(p_{n+1}\in D_n\cap N\) with \(p_{n+1}\leq_P p_n\).
A fixed enumeration of \(N\) selects its first suitable witness at each
step, making a specified ordinary recursive construction.
It is not claimed to be a computable procedure for arbitrary set-theoretic data.

Countable closure in the metatheory gives \(q\leq_P p_n\) for every \(n\).
For any \(r\leq_P q\), \(r\leq_P p_{n+1}\in D_n\cap N\),
so \(r\) itself witnesses compatibility with that member.
This proves predensity for every dense set in the enumeration and
therefore the master property. In particular \(q\leq_P p\).
Neither \(q\in N\) nor \(q\in D_n\) was assumed. If \(D_n\) is open,
the latter membership does follow, but that is not needed for general \(D_n\).
\end{proof}
\end{annexitem}

\begin{annexitem}{annex:p06.pfm.examples}{thm:p06-pfm-examples}{Cohen strings and countable partial functions}
\begin{proof}
Enumerating binary strings by length and then lexicographically proves
\(2^{<\omega}\) countable. Its empty string is greatest, so the countable
properness theorem applies. The descending chain of strings consisting
of \(n\) zeroes has no finite-string lower bound: such a bound would
have length at least every \(n\). Hence the poset is not countably closed.

For a ground-model generic filter, compatible strings are comparable,
and their union is a binary function. The dense sets of strings of length
at least \(n\) make its domain all of \(\omega\).
For each ground-model binary real \(x\), strings which disagree with \(x\)
are dense: any string already disagreeing qualifies; otherwise extend it
by the bit opposite \(x\) at its next position.
The generic meets each such dense set, so its union differs from every
ground-model real. This assertion is conditional on the generic extension,
not the existence of a generic filter in the ground universe itself.

Now let \(Q\) consist of countable partial maps \(\omega_1\to2\).
A descending chain is an increasing chain of compatible function graphs.
Its union is a function, and its domain is countable: Choice supplies
enumerations of the nonempty domains; pairing the enumeration index with
the entry index enumerates their union. The empty domains contribute nothing.
Thus the union is a condition and a common lower bound.
Countable closure proves properness.
For each \(x\in2^\omega\), the condition with domain \(\omega\) and values
\(x\) belongs to \(Q\). Distinct such conditions disagree somewhere and
are incompatible. They form an uncountable antichain: for any enumeration
of binary reals, the diagonal opposite-bit real is omitted.
Thus \(Q\) is not ccc. Both examples have a greatest condition and meet
the same declared height convention.
\end{proof}
\end{annexitem}

\begin{annexitem}{annex:p06.pfm.semigeneric}{thm:p06-pfm-semigeneric}{Ordinal-name decisions on a club of models}
\begin{proof}
We state the imported forcing interface explicitly.
For a name \(\tau\) with \(1_P\forces\tau<\check\omega_1\),
conditions deciding \(\tau\) as a ground ordinal form a dense set.
Monotonicity preserves a decided value under strengthening.
If conditions forcing a fixed assertion are dense below \(q\), then
\(q\) forces that assertion. These are the ordinal-decision and
truth/density consequences of the cited recursive forcing theorem,
not newly assumed properties of an arbitrary binary relation.

Fix a height and let \(T_\chi\) be the set of relevant names in \(H_\chi\).
Choose for each \(\tau\in T_\chi\) a maximal deciding antichain
\(A_\tau\subseteq P\), using the proved dense-antichain construction.
Let \(f_\tau:A_\tau\to\omega_1\) give its decided ordinal values.
Both \(A_\tau\) and \(f_\tau\) belong to \(H_\chi\):
their transitive closures are bounded by that of the poset datum,
\(\omega_1\), and the finite pair coding, of size at most the declared
\(\kappa<\chi\).
Use Choice over the set \(T_\chi\) to select these pairs.
Expand \((H_\chi,\in)\) by the two unary operations returning these sets
and by a predicate for \(T_\chi\), using fixed defaults outside it.
This is a finite, hence countable, language expansion.
The hull theorem supplies a club of countable models elementary in it.
Their membership reducts are elementary and contain the named poset datum.

Let \(N\) be one of these models, \(q\) a master condition, and
\(\tau\in T_\chi\cap N\).
The recorded \(A_\tau,f_\tau\) belong to \(N\).
The dense \(A_\tau^\downarrow\) is uniquely defined from that antichain
and the order, so also belongs to \(N\).
For each \(r\leq_P q\), the master test supplies
\(d\in A_\tau^\downarrow\cap N\) compatible with \(r\).
Elementarity supplies \(a\in A_\tau\cap N\) with \(d\leq_P a\).
Unique function evaluation puts \(f_\tau(a)\) in \(N\cap\omega_1=\delta_N\).
A common strengthening \(s\leq_P r,d\) is below \(a\), and hence forces
\(\tau=\check{f_\tau(a)}<\check{\delta_N}\).
Such conditions are dense below \(q\), so the imported density clause
gives \(q\forces\tau<\check{\delta_N}\).
This holds for every relevant name in \(N\), proving semigenericity.
Properness supplies master strengthenings for every \(p\in P\cap N\).
The chosen club therefore witnesses the local semiproperness test.
Repeat at every admissible height.
The deciding operations were added explicitly; no absoluteness of the
forcing predicate to \(H_\chi\) was used.
\end{proof}
\end{annexitem}

\begin{annexitem}{annex:p06.pfm.collapse}{thm:p06-pfm-collapse}{The extension which destroys the master test}
\begin{proof}
The finite-map poset has greatest member empty function.
For every natural \(n\), \(D_n=\{s:n\in\operatorname{dom}(s)\}\)
is dense: add the pair \((n,0)\) if needed.
It is a subset of \(P\), belongs to \(H_\chi\), and is uniquely defined
there from \(n\) and the poset datum.
Since \(n\in N\) and \(\mathbf P\in N\), it belongs to \(N\).

Fix arbitrary \(q\) and choose \(n\) outside its finite domain.
Since \(\delta_N<\omega_1\), \(r=q\cup\{(n,\delta_N)\}\) is a condition
stronger than \(q\).
For \(s\in D_n\cap N\), unique function evaluation at \(n\) gives
\(s(n)\in N\cap\omega_1\), so \(s(n)<\delta_N\).
The two functions assign different values at \(n\); no common extension
exists. Thus \(D_n\cap N\) is not predense below \(q\).
This holds for every \(q\), including any candidate below \(p\in P\cap N\),
and proves failure of properness.

For the semigeneric claim we identify the names, not just assert that
a collapse must fail. In a transitive ground model let
\[
 \tau_n=\{(\check\gamma,p):p\in P,\ n\in\operatorname{dom}(p),\
                                      \gamma<p(n)\}.
\]
The standard generic evaluation and compatibility of finite functions
give \(\tau_n^G=g(n)\), where \(g=\bigcup G\).
Indeed, if \(p\in G\) decides the \(n\)-th value as \(\alpha\),
the members of this interpreted name are exactly the ordinals
\(\gamma<\alpha\); every other member of \(G\) assigning an \(n\)-th
value assigns that same \(\alpha\).
Thus \(1_P\forces\tau_n<\check\omega_1\), and an assigning condition
forces \(\tau_n\) equal to its assigned ground ordinal.

All these names and their sequence belong to \(H_\chi\).
To make the elementarity step independent of any internal forcing
absoluteness, include the explicit ledger
\(\{(\gamma,\check\gamma):\gamma<\omega_1\}\) and the sequence
\(\langle\tau_n:n<\omega\rangle\) as named parameters.
Their hereditary sizes are at most \(\kappa\); the canonical-name
recursion adds only the finite pair/poset coding to the \(\omega_1\)
ordinal ledger. The hull theorem gives a club of models containing them.
For each such \(N\), unique evaluation puts every \(\tau_n\) in \(N\).
The \(r\) just constructed forces \(\tau_n=\check{\delta_N}\);
therefore its weaker \(q\) cannot force
\(\tau_n<\check{\delta_N}\), by monotonicity and consistency of the
forcing truth clauses. No \(q\) is semigeneric for these models.
If a club witnessed semiproperness, intersect it with this ledger club.
The proved club-intersection lemma gives a model in both, where every
condition would have a semigeneric strengthening and none can have one.
This contradiction proves failure of the club test.

Finally, for every \(\alpha<\omega_1\),
\(E_\alpha=\{p:\alpha\in\operatorname{ran}(p)\}\) is dense: assign
\(\alpha\) at a fresh natural coordinate.
A generic filter meets all \(D_n\) and \(E_\alpha\).
Its directedness makes its union a function; the first family gives
domain \(\omega\), and the second gives range the whole ground \(\omega_1\).
This proves the claimed collapse without assuming a generic exists
inside its ground model.
\end{proof}
\end{annexitem}
